\section{Data for additional volume and surface averaged tallies}
\label{sec2.10}

\medskip

\textbf{General remarks}

\medskip

There is a large number of preprogrammed ``default" volume or
surface averaged tallies, which have been selected mainly to allow
assessment of global particle and energy balances for all test
particle species, as well as coupling of neutral gas and plasma
transport equations. These tallies are estimated at each EIRENE
run unless they are explicitly abandoned by the surface crossing
switches ILSWCH (\cref{sec2.3b}), or are turned off in
input block 11 (\cref{sec2.11}).

There is, however, a very general option to allow for estimation
of many more moments of the test particle $\mu$ -space
distribution function by resorting to appropriately prepared
problem specific routines UPTUSR, UPCUSR and UPSUSR. These are
described in more detail in \cref{sec3.2}. Furthermore
there are options to form algebraic expressions from tallies, and,
since Eirene$_{2003}$ also for particle and energy flux spectra
evaluated from the fluxes on selected surfaces or in selected
cells of the computational grid. The input data in block 10
described here are flags for activating and controlling these
routines.

Default: If the first card in this block: \lstinline|NADVI|, \dots,
contains only zeros, then the rest of this block can/must be
omitted.

\medskip

\textbf{The Input Block}

\begin{lstlisting}
*** 10. Data for additional tallies
      NADVI  NCLVI  NALVI  NADSI  NALSI NADSPC
\end{lstlisting}
\textsl{* 10A. Data for additional track-length estimated volumetric
tallies}
\begin{lstlisting}
      DO 101 IADVI=1,NADVI
        IADVE(IADVI) IADVS(IADVI) IADVT(IADVI) IADVR(IADVI)
        TXTTAL(IADVI,NTALA)
        TXTSPC(IADVI,NTALA) TXTUNT(IADVI,NTALA)
  101 CONTINUE
\end{lstlisting}
\textsl{* 10B. Data for additional collision estimated volumetric
            tallies}
\begin{lstlisting}
      DO 102 ICLVI=1,NCLVI
        ICLVE(ICLVI) ICLVS(ICLVI) ICLVT(ICLVI) ICLVR(ICLVI)
        TXTTAL(ICLVI,NTALC)
        TXTSPC(ICLVI,NTALC) TXTUNT(ICLVI,NTALC)
  102 CONTINUE
\end{lstlisting}
\textsl{* 10C. Data for algebraic function of volumetric
             tallies}
\begin{lstlisting}
      DO 103 IALVI=1,NALVI
        ALSTRNG
        TXTTAL(IALVI,NTALR)
        TXTSPC(IALVI,NTALR) TXTUNT(IALVI,NTALR)
  103 CONTINUE
\end{lstlisting}
\textsl{* 10D. Data for additional surface crossing tallies}
\begin{lstlisting}
      DO 104 IADSI=1,NADSI
        IADSE(IADSI) IADSS(IADSI) IADST(IADSI) IADSR(IADSI)
        TXTTLS(IADSI,NTLSA)
        TXTSPS(IADSI,NTLSA) TXTUNS(IADSI,NTLSA)
  104 CONTINUE
\end{lstlisting}

\textsl{* 10E. Data for algebraic function of
surface crossing tallies}
\begin{lstlisting}
      DO 105 IALSI=1,NALSI
        ALSTRNG
        TXTTLS(IALSI,NTLSR)
        TXTSPS(IALSI,NTLSR) TXTUNS(IALSI,NTLSR)
  105 CONTINUE
\end{lstlisting}

optional, in versions Eirene$_{2003}$ and younger:

\textsl{* 10F. Data for energy spectra}
\begin{lstlisting}
      DO 106 IADSP=1,NADSPC
        ISPSRF IPTYP IPSPZ ISPTYP NSPS ISRFCLL IDIREC
        SPCMN SPCMX SPC_SHIFT SPCPLT_X SPCPLT_Y SPCPLT_SAME
        if (idirec.ne.0) SPCVX SPCVY SPCVZ
  106 CONTINUE
\end{lstlisting}

\textbf{Meaning of the input variables for additional volume and
surface averaged Tallies}

\medskip

\begin{description}
\item[NADVI] Total number of additional volume averaged,
track-length estimated tallies
\item[NCLVI] Total number of
additional volume averaged, collision estimated tallies
\item[NALVI] Total number of tallies defined as algebraic expressions
of other volume averaged tallies
\item[NADSI] Total number
of additional surface averaged tallies
\item[NALSI] Total
number of tallies defined as algebraic expressions of other
surface averaged tallies
\item[NADSPC] Total number of
surface or cell averaged energy spectra
\end{description}
\subsection{Additional volume averaged tallies, track-length estimator}\label{sec2.10A}
\begin{description}
\item[IADVE] flag
for scaling factor for this tally (carried out in subroutine
MCARLO)
\begin{description}
\item[$= 1$]
scale tally per unit volume [\si{\per\cubic\centi\metre}]. The tally is printed and
plotted in the units

$[\si{\per\cubic\centi\metre}] \cdot [$units of $g^*] \cdot [\si{\centi\metre}] \cdot
[$source strength FLUX$]$

however, with FLUX converted to units [\si{\per\second}] (rather than input units $[Ampere]$).
For the definition of the detector functions
$g$ and $g^*$ see \cref{sec3.2}, the variable FLUX is explained in
input block no.\ 7 (\cref{sec2.7}). This scaling is default for ``density tallies" of particles, momentum and energy.
\item[$= 2$]
scale tally per unit cell. Same units as above, however not per
\si{\cubic\centi\metre} but per cell instead.
\item[$= 3$]
same as IADVE = 1, but with FLUX in Ampere, rather than \si{\per\second}.
This scaling is default for ``source rate tallies", e.g.\ for particle, momentum and energy sources.
\item[$= 4$]
same as IADVE = 2, but with FLUX in Ampere, rather than \si{\per\second}.
\item[else]
no re-scaling done, units as chosen in subroutine UPTUSR.
\end{description}
\end{description}
Note: scaling of EIRENE default tallies 1 to 8 (particle and energy
densities, respectively) and tallies 85 to 96 (momentum density) is as for IADVE=1 option. Scaling of all
default source and sink terms is as for IADVE=3 option.

If instead also for particle- or energy densities tallies 1 to 8 the scaling IADVE=3 would
have been used, then, for example, particle densities would have
been in units [\si{\ampere\second\per\cubic\centi\metre}] (``atomic charge density")
rather than [\si{\per\cubic\centi\metre}]. Energy densities would have been in $[\si{\watt}
\cdot \si{\second\per\cubic\centi\metre}]$ rather than $[\si{\electronvolt\per\cubic\centi\metre}]$.
\begin{description}
\item[IADVS]
species index of default volume averaged tally,
which is to be replaced by this additional track-length estimated tally.
In case of tallies with no species index, IADVS must be set equal to 1.
\item[IADVT]
number of default volume averaged tally,
which is to be replaced by this track-length estimated tally.
\end{description}
If either IADVS or IADVT is out of range, no default tally is
replaced by this additional tally.
\begin{description}
\item[IADVR]
If automatic re-scaling of volumetric tallies is performed (i.e.
if NLSCL = TRUE), then re-scale this tally with EIRENE recommended
factor FATM (IADVR = 1), FMOL (IADVR = 2) or FION (IADVR = 3) (see
block 1 for the flag NLSCL and the variables FATM,FMOL,FION)
\item[TXTTAL]
Text to label tally on numerical or graphical output.
\item[TXTSPC]
Text describing the species of particles contributing for
this tally in output routines.
\item[TXTUNT]
Text describing the units of this tally in output routines.
\end{description}
\subsection{Additional volume averaged tallies, collision estimator}\label{sec2.10B}
\begin{description}
\item[ICLVE]
as IADVE above, for collision estimated tally (subroutine UPCUSR).
\item[ICLVS]
species index of default volume averaged tally,
which is to be replaced by this collision estimated tally.
In case of tallies with no species index, ICLVS must be set equal to 1.
\item[ICLVT]
number of default volume averaged tally,
which is to be replaced by this collision estimated tally.
\end{description}
If either ICLVS or ICLVT is out of range, no default tally is
replaced.
\begin{description}
\item[ICLVR] as IADVR above, for collision estimated tally
(subroutine UPCUSR).
\item[TXTTAL, TXTSPC, TXTUNT]: as above
\end{description}
\subsection{Algebraic expression in volume averaged tallies}\label{sec2.10C}
\begin{description}
\item[ALSTRNG] character string which is interpreted as an
algebraic expression in some volume averaged tallies. An operand
$\langle i,j \rangle$ stands for tally number $j$, first (species) index $i$, in
\cref{tab1} (input tallies) and \cref{tab2} (output tallies) . Note that the first index for
tallies with no species index must read 1, and the tally number of input tallies must be negative.
Example: $\langle 1,-1\rangle$ for the electron temperature tally, $\langle 2,3\rangle$ for the particle density of test ion species
no.\ IION=2 . Expressions $\langle c\rangle$
with an integer or real constant $c$ are interpreted as scalars. The
string may contain an arbitrary (but $\le$ 20) number of operands,
and of operators +, -, *, /, **, and of properly nested parentheses
(\dots).

Standard deviations are not available, generally, for algebraic tallies.
For linear combinations of tallies it is, in principle, possible to obtain also the standard deviations,
this evaluation of error estimates for such combinations of tally is, however, is currently  carried
out only in a proprietary code segment (available from the author).

E.g., the total electron particle source due to test particle - plasma
interaction in units: \si{\#\per\second\per\cubic\metre} can be obtained by the line:

($\langle 1,7\rangle +\langle 1,12\rangle +\langle 1,17\rangle) * \langle 1.e6\rangle /\langle 1.6022e-19\rangle$

in code versions older than 2002 (see tables in
\cref{sec.tables-old}), and the same expressions in versions 2002
and younger, i.e.\ after implementation of photons as further
species type (and the related default tallies, tables in
\cref{sec.tables-new}):

($\langle 1,9\rangle+\langle 1,15\rangle+\langle 1,21\rangle) * \langle 1.e6\rangle/\langle 1.6022e-19\rangle$

and it would be stored on the tally: ``ALGV" with the first
(labelling) index IALVI.

Note that the cell volume array VOL is regarded as a volume averaged
tally, by abuse of language (tally number = -14, see \cref{tab1}).
\item[TXTTAL, TXTSPC, TXTUNT]
as above.
\end{description}
\subsubsection{frequently used algebraic tallies}
\begin{itemize}
\item{mean energy} ~

The mean energy (\si{\electronvolt}) of a test particle species is the ratio of its total energy density (= total pressure) and particle density.
The algebraic tally
$\langle i,5\rangle/\langle i,1\rangle$ is that ratio estimate, here exemplarily for atomic species no.\ IATM=$i$.
Similarly, for molecular species IMOL=$k$, this expression would
read: $\langle k,6\rangle/\langle k,2\rangle$

\item{kinetic temperature} ~

The kinetic temperature $T_t$ (\si{\electronvolt})  is defined such that $3/2  \ T_t$ is the mean thermal energy. The kinetic temperature
is obtained by first subtracting the dynamic pressure (contribution from drift motion) from the total energy density,
and then by rescaling.
\begin{align} \nonumber
(\langle i,5\rangle&-\langle 1.8792e+35\rangle/\langle 2.0\rangle \\
&*(\langle i,85 \rangle*\langle i,85\rangle+\langle i,89\rangle*\langle i,89\rangle+\langle i,93\rangle*\langle i,93\rangle)\\
&/\langle i,1\rangle)/\langle i,1\rangle/\langle 1.5\rangle
\end{align}
Here the subtracted part
$\langle 1.8792e+35\rangle/\langle 2.0\rangle*
(\langle i,85\rangle*\langle i,85\rangle+
 \langle i,89\rangle*\langle i,89\rangle+
 \langle i,93\rangle*\langle i,93\rangle)/
 \langle i,1\rangle)$ is the dynamic pressure, scaled to \si{\electronvolt\cubic\centi\metre} units
of atoms, with atomic weight $\langle 2.0\rangle$ in atomic mass units in this example.
The numerical factor $1.8792e+35$ is, in EIRENE unit conversion variables, is given
by CVELI2/AMUA/AMUA, (defined in SETCON.F), and the momentum density tallies 85, 89 and 93 are in \si{\gram\centi\metre\per\second\per\cubic\centi\metre}.
\end{itemize}
\subsection{Additional surface averaged tallies}\label{sec2.10E}

as in sub-block 10A or 10B, but for surface averaged tallies rather than
volume averaged tallies.


Note: for surface averaged tallies collision and track-length
estimators are identical, see \cref{sec3.2.4}.

\subsection{Algebraic expression in surface averaged tallies}\label{sec2.10F}



as in sub-block 10C, but for surface averaged tallies rather
than volume averaged tallies. The corresponding tally numbers are listed
in \cref{tab3}.

There is one generalization available here compared to the options
for volume tallies, which permits algebraic expressions not only
between different complete surface tallies (expression evaluated
for all NLIMPS surfaces), but also between different individual
surfaces.

If the first index I is out of range, i.e., larger than the
first dimension N1 of the tally J as listed in \cref{tab3},
then the operand $\langle I,J\rangle$ identifies the following estimate:

surface index $ILIMPS = (I-N1-1)/N1 + 1 $

species index $ISPZ = I - ILIMPS \cdot N1 $

The other way round: in order to identify species ISPZ and surface no.\ ILIMPS
for a particular tally J,
on must set:

first operand $I = ILIMPS \cdot N1 + ISPZ  $

second operand  J

The algebraic expression is stored on the ALGS tally for all
surfaces ILIMPS, which occurred in at least one of the operands.

Example: let there be up to 3 atom species (N1=NATM=3), and assume, that
additional surfaces no.\ 7,8,9,10 and 11 comprise the pump.
Let furthermore IATM=2 stand for helium atoms.
The total pumped flux of helium atoms from these surfaces is
the difference between incident and re-emitted He particle fluxes,
summed over these 5 surfaces. This is evaluated
by the string:

$\langle 23,1\rangle-\langle 23,2\rangle+\langle 26,1\rangle-\langle 26,2\rangle+\langle 29,1\rangle-$

$\langle 29,2\rangle+\langle 32,1\rangle-\langle 32,2\rangle+\langle 35,1\rangle-\langle 35,2\rangle$

The value found for this expression is stored on ALGS (\dots, ISURF),
and is the same for ISURF=7, =8, =9, =10 or =11, and zero for all
other surfaces.

\subsection{Energy spectra in selected cells or surfaces}\label{sec2.10G}
Additional, spectrally resolved, tallies, either surface averaged or cell based.
Printout of these spectra is activated in input block 11a, flag NSPCPR.

\begin{description}
\item[ISRFCLL] flag to choose between surface and cell averages.

=0:  surface averaged spectrum

=1:  volume (cell) averaged spectrum
\item[ISPSRF]
number of surface or cell, for which spectra are to be evaluated.
Positive for additional surfaces, negative for non-default standard
surfaces.
\item[IPTYP]
Type of particle (0 for photons, 1 for atoms, 2 for molecules,
3 for test ions, 4 for bulk ions)
\item[IPSPZ]
species index of particle, within the specified type-category.
If IPSPZ=0, then sum over species, for this type of particles.


\item[ISPTYP]
type of spectrum:

\textbf{for surface averaged spectra (i.e.\ ISRFCLL=0):}

flux-spectrum (ISPTYP=1) , [\si{\ampere\per\electronvolt}]

or energy-flux spectrum (ISPTYP=2) [\si{\watt\per\electronvolt}]

Currently only surface ``incident fluxes". Spectra for surfaces ``emitted fluxes" (outgoing)
are not yet available.

\textbf{for cell averaged spectra (i.e.\ ISRFCLL=1):}

\textbf{case a:} no direction specified: IDIREC=0

energy distribution (ISPTYP=1) , [\si{\per\cubic\centi\metre\per\electronvolt}]

energy weighted energy distribution (ISPTYP=2) [\si{\electronvolt\per\cubic\centi\metre\per\electronvolt}]

velocity weighted energy distribution (ISPTYP=3) [\si{\centi\metre\per\second\per\cubic\centi\metre\per\electronvolt}]

\textbf{case b:} direction specified: IDIREC=1


velocity distribution along specified direction, binned in energy units (ISPTYP=1), [\si{\per\cubic\centi\metre\per\electronvolt}]

energy weighted velocity distribution along specified direction, put into bins in energy units
(ISPTYP=2) [\si{\electronvolt\per\cubic\centi\metre\per\electronvolt}]

velocity weighted velocity distribution along specified direction,
put into bins in energy units (ISPTYP=3) [\si{\centi\metre\per\second\per\cubic\centi\metre\per\electronvolt}]


\item[IDIREC, SPCVX, SPCVY, SPCVZ]

These flags are only for cell averaged spectra (otherwise: not in use):

=0: no direction specified

=1: cell averaged spectra are along direction SPCVX, SPCVY, SPCVZ. The vector
defining this direction is normalized internally.

\item[NSPS, SPCMN, SPCMX] Number of equidistant energy
bins (\si{\electronvolt}) for spectra. The contributions with energies below the
minimum energy SPCMN are stored in bin no.\ 0, those with energies
above the maximum energy SPCMX are stored in bin NSPS+1.

NSPS $> 0$:  The bins are set equidistant in linear energy scale

NSPS $< 0$:  $\mid$NSPS$\mid$ is used for the number of bins, but the bins are set equidistant
in a logarithmic energy scale (option added by X.B., 2017). Same as for NCHENI parameter in input
block 12 for (post processed) line of sight integrated and spectrally resolved data.

\item[SPC\_SHIFT] shift spectrum by SPC\_SHIFT (\si{\electronvolt}) in printout and plots.

\item[SPCPLT\_X, SPCPLT\_Y] SPCPLT\_X is a flag foreseen to control the
X axis (independent variable) in a plot of the spectrum. (X: e.g.\ the no.\ of an energy bin).
This flag is currently not in use.

SPCPLT\_Y is a flag to control the
Y axis (dependent variable) in a plot of the spectrum. (Y: e.g.\ particle or energy flux in an energy bin).
IF SPCPLT\_Y $>$ 0: logarithmic scale for Y axis, otherwise: linear scale for Y axis.
.

\item[SPCPLT\_SAME] plot this spectrum into previous
frame
\end{description}
The surface flux spectra are printed together with the other
surface averaged tallies, for those surfaces, which are selected
for printout in sub-block 11A. They are all automatically plotted
as soon as at least one further volume tally plot (input sub-block
11B.2) is selected.
